\documentclass[11pt]{article}

\usepackage[margin=0.92in]{geometry}
\usepackage{amsmath,amssymb,amsthm,mathtools}
\usepackage{booktabs}
\usepackage{microtype}
\usepackage{enumitem}
\usepackage[hidelinks]{hyperref}

\newtheorem{theorem}{Theorem}[section]
\newtheorem{lemma}[theorem]{Lemma}
\newtheorem{proposition}[theorem]{Proposition}
\newtheorem{corollary}[theorem]{Corollary}
\newtheorem{definition}[theorem]{Definition}
\theoremstyle{remark}
\newtheorem{remark}[theorem]{Remark}

\newcommand{\Pacc}{\widehat P}
\newcommand{\RR}{\mathbb{R}}

\title{\textbf{Scale-Optimized Geometric Acceleration of Polygonal Approximations to $\pi$}}
\author{C. Wayne Baker\\Independent Researcher}
\date{May 9, 2026; revised July 2026}

\begin{document}
\maketitle

\begin{abstract}
Let
\[
  P_N=N\sin\!\left(\frac{\pi}{N}\right)
\]
be the half-perimeter of the regular $2N$-gon inscribed in the unit circle.
Its error contains only even powers of $N^{-1}$.  For an integer scale
factor $b\ge2$ and an order $r\ge1$, we form
\[
  \Pacc^{[r]}_{N,b}
  =\sum_{m=0}^{r-1}w_m^{(r,b)}P_{b^mN},
\]
where the weights are the Lagrange extrapolation weights at $x=0$ for the
geometric nodes $x_m=b^{-2m}$.  The first $r-1$ even-power errors cancel,
so
\[
  \Pacc^{[r]}_{N,b}-\pi=O(N^{-2r}).
\]
For geometric nodes, the weights are uniformly stable in $r$:
\[
  \sum_{m=0}^{r-1}|w_m^{(r,b)}|
  \le
  S_b,
  \qquad
  S_b=\prod_{\ell=1}^{\infty}
  \frac{1+b^{-2\ell}}{1-b^{-2\ell}}.
\]
An explicit sine-remainder bound then quantifies the accelerated error.
High-precision computations verify the predicted slopes through $r=24$ for
$b=2,3,6,9,27$.

The constructibility status is separated carefully.  The dyadic branch is
exactly compass-and-straightedge constructible by nested square roots.  Exact
tripling requires a cubic inversion and is not generally constructible.  A
constructible linear chord seed has error $O(N^{-3})$, and a general
perturbation theorem states the accuracy required of any approximate tripling
mechanism to preserve an $r$-scale hierarchy.  No exact construction of every
regular $3N$-gon is claimed.
\end{abstract}

\section{Introduction and Claim Boundary}

The classical polygon sequence
\[
  P_N=N\sin\!\left(\frac{\pi}{N}\right)
\]
converges to $\pi$ with second-order error.  The important fact is not only
that the error is small, but that it is structured:
\[
  P_N
  =\pi+\frac{C_2}{N^2}+\frac{C_4}{N^4}+\frac{C_6}{N^6}+\cdots.
\]
This permits exact cross-scale cancellation of successive error terms.

The paper has four principal claims.
\begin{enumerate}[label=(\roman*)]
  \item Geometric scale weights cancel the first $r-1$ even-power modes.
  \item For nodes $b^{-2m}$, the total weight amplification is uniformly
        bounded in $r$.
  \item The remaining analytic error has an explicit factorial bound.
  \item Scale selection is a tradeoff: increasing $b$ improves weight
        stability but increases geometric work.
\end{enumerate}

The word \emph{geometric} refers to the polygonal source sequence and its
scale hierarchy.  It does not imply that every scale branch is exactly
straightedge-and-compass constructible.  The branch $b=2$ is exact and
constructible.  Branches containing a factor of $3$ are analytic or numerical
scale architectures unless a separate approximate trisection construction is
supplied with a sufficient error certificate.

\section{The Even-Power Polygonal Expansion}

Using the sine series,
\[
  \sin z=z-\frac{z^3}{3!}+\frac{z^5}{5!}-\frac{z^7}{7!}+\cdots,
\]
with $z=\pi/N$, one obtains
\[
  P_N
  =\pi-\frac{\pi^3}{6N^2}
  +\frac{\pi^5}{120N^4}
  -\frac{\pi^7}{5040N^6}+\cdots.
\]
More generally,
\[
  P_N
  =\pi+\sum_{j=1}^{\infty}
  (-1)^j\frac{\pi^{2j+1}}{(2j+1)!}N^{-2j}.
\]
Only even powers of $N^{-1}$ occur because $\sin z/z$ is an even analytic
function.

\begin{definition}[Even-power error structure]
A sequence $Q_N\to Q$ has an even-power error structure if
\[
  Q_N=Q+C_2N^{-2}+C_4N^{-4}+C_6N^{-6}+\cdots.
\]
\end{definition}

\section{Scale-$b$ Acceleration}

Fix an integer $b\ge2$ and an order $r\ge1$.  Define
\[
  x_m=b^{-2m},\qquad m=0,1,\ldots,r-1.
\]
The Lagrange weights for evaluating at $x=0$ are
\[
  w_m^{(r,b)}
  =\prod_{\substack{0\le i\le r-1\\ i\ne m}}
  \frac{-x_i}{x_m-x_i}.
\]
The accelerated polygonal approximation is
\[
  \Pacc^{[r]}_{N,b}
  =\sum_{m=0}^{r-1}w_m^{(r,b)}P_{b^mN}.
\]
The weights satisfy
\[
  \sum_{m=0}^{r-1}w_m^{(r,b)}=1
\]
and, for $j=1,\ldots,r-1$,
\[
  \sum_{m=0}^{r-1}w_m^{(r,b)}b^{-2jm}=0.
\]

\begin{theorem}[Scale-$b$ cancellation]
Suppose
\[
  Q_N=Q+\sum_{j=1}^{r-1}C_{2j}N^{-2j}+O(N^{-2r}).
\]
Then
\[
  \widehat Q^{[r]}_{N,b}
  :=\sum_{m=0}^{r-1}w_m^{(r,b)}Q_{b^mN}
  =Q+O(N^{-2r}).
\]
\end{theorem}

\begin{proof}
Substitution gives
\[
\widehat Q^{[r]}_{N,b}
=Q\sum_mw_m^{(r,b)}
+\sum_{j=1}^{r-1}C_{2j}N^{-2j}
 \sum_mw_m^{(r,b)}b^{-2jm}
+O(N^{-2r}).
\]
The unity condition preserves $Q$, and the moment conditions cancel the first
$r-1$ error modes.
\end{proof}

For $r=2$,
\[
  \Pacc^{[2]}_{N,b}
  =P_{bN}+\frac{P_{bN}-P_N}{b^2-1}
  =\pi+O(N^{-4}).
\]
For $b=2$ and $r=3$,
\[
  \Pacc^{[3]}_{N,2}
  =\frac{1}{45}P_N-\frac{4}{9}P_{2N}
   +\frac{64}{45}P_{4N}
  =\pi+O(N^{-6}).
\]
The coefficient $1/45$ is essential; a surviving draft incorrectly displayed
$1/64$.

\section{Geometric-Node Weight Stability}

Set $q=b^{-2}$ and use the notation
\[
  (q;q)_n=\prod_{\ell=1}^{n}(1-q^\ell),
  \qquad (q;q)_0=1.
\]

\begin{lemma}[Closed weight formula]
Let $s=r-1-m$.  Then
\[
  |w_m^{(r,b)}|
  =\frac{q^{s(s+1)/2}}{(q;q)_m(q;q)_s}.
\]
\end{lemma}

\begin{proof}
Split the defining product into factors with $i<m$ and $i>m$.  After
factoring the powers of $q$, the remaining products are precisely
$(q;q)_m$ and $(q;q)_s$.
\end{proof}

\begin{theorem}[Uniform stability]
For every $r\ge1$,
\[
  \sum_{m=0}^{r-1}|w_m^{(r,b)}|
  \le S_b,
\]
where
\[
  S_b
  =\prod_{\ell=1}^{\infty}
   \frac{1+b^{-2\ell}}{1-b^{-2\ell}}<\infty.
\]
\end{theorem}

\begin{proof}
Writing $s=r-1-m$ and using $(q;q)_{r-1-s}\ge(q;q)_\infty$,
\[
\sum_{m=0}^{r-1}|w_m^{(r,b)}|
\le\frac{1}{(q;q)_\infty}
\sum_{s=0}^{\infty}\frac{q^{s(s+1)/2}}{(q;q)_s}.
\]
Euler's product identity gives
\[
  \sum_{s=0}^{\infty}\frac{q^{s(s+1)/2}}{(q;q)_s}
  =(-q;q)_\infty.
\]
Therefore the bound is
\[
  \frac{(-q;q)_\infty}{(q;q)_\infty}
  =\prod_{\ell=1}^{\infty}\frac{1+q^\ell}{1-q^\ell}.
\]
\end{proof}

The corrected constants for the tested branches are shown in
Table~\ref{tab:stability}.

\begin{table}[ht]
\centering
\begin{tabular}{ccc}
\toprule
$b$ & $q=b^{-2}$ & $S_b$\\
\midrule
2  & $1/4$   & $1.969260353668269\ldots$\\
3  & $1/9$   & $1.285210586391831\ldots$\\
6  & $1/36$  & $1.058822194456829\ldots$\\
9  & $1/81$  & $1.025316406852818\ldots$\\
27 & $1/729$ & $1.002751031629743\ldots$\\
\bottomrule
\end{tabular}
\caption{Uniform stability constants.  The May concept PDF rounded the
$b=2$ constant incorrectly; the value above is independently recomputed.}
\label{tab:stability}
\end{table}

\section{Explicit Analytic Remainder Bound}

\begin{theorem}[Sine-remainder bound]
Let $N\ge6$, $b\ge2$, and $r\ge1$.  Then
\[
  \left|\Pacc^{[r]}_{N,b}-\pi\right|
  \le
  \frac{\pi^{2r+1}}{(2r+1)!}N^{-2r}
  \sum_{m=0}^{r-1}|w_m^{(r,b)}|b^{-2rm}.
\]
Consequently,
\[
  \left|\Pacc^{[r]}_{N,b}-\pi\right|
  \le
  S_b\frac{\pi^{2r+1}}{(2r+1)!}N^{-2r}.
\]
\end{theorem}

\begin{proof}
At scale $b^mN$, subtract the first $r-1$ nonconstant sine terms.  Since
$\pi/(b^mN)\le\pi/6<1$, the alternating Taylor remainder is bounded by its
first omitted term:
\[
  |R_r(b^mN)|
  \le\frac{\pi^{2r+1}}{(2r+1)!}(b^mN)^{-2r}.
\]
The acceleration weights cancel the retained terms.  Applying the triangle
inequality yields the first estimate; dropping $b^{-2rm}\le1$ and invoking
the stability theorem gives the second.
\end{proof}

\section{High-Precision Verification}

The accompanying script evaluates the weights, moment residuals, stability
constants, accelerated errors, and observed slopes.  For a fixed $b$ and $r$,
the slope is obtained by a least-squares fit of
\[
  \log_{10}|\Pacc^{[r]}_{N,b}-\pi|
\]
against $\log_{10}N$ over $N=16,20,24,32$.  The predicted slope is $-2r$.

\begin{table}[ht]
\centering
\begin{tabular}{ccc}
\toprule
Scale $b$ & Order $r$ & Observed slope\\
\midrule
2  & 24 & $-47.9999784973\ldots$\\
3  & 24 & $-47.9999818571\ldots$\\
6  & 24 & $-47.9999834123\ldots$\\
9  & 24 & $-47.9999836714\ldots$\\
27 & 24 & $-47.9999838509\ldots$\\
\bottomrule
\end{tabular}
\caption{Representative 1300-digit slope audit.  The predicted value is
$-48$.}
\label{tab:slopes}
\end{table}

These computations are checks of the formulas, not substitutes for the
proofs.  Their main use is to confirm that high-order cancellation remains
numerically visible when sufficient working precision is used.

\section{Scale Choice: Stability Versus Work}

The stability constant satisfies
\[
  S_b\longrightarrow1\qquad(b\to\infty),
\]
so larger scale ratios reduce coefficient amplification.  This does not by
itself identify an optimal scale.  If the cost of constructing or evaluating a
polygon is proportional to its side count, a simple work proxy is
\[
  W_b(r)=\sum_{m=0}^{r-1}b^m=\frac{b^r-1}{b-1}.
\]
For fixed $r$, $W_b(r)$ grows rapidly with $b$, while $S_b$ improves only
toward $1$.  Thus scale selection balances at least three quantities:
\[
  \text{scale separation},\qquad
  \text{weight stability},\qquad
  \text{construction cost}.
\]
There is no universal optimum without a specified cost model.  Earlier claims
that one tested branch is globally ``best'' should be read only relative to
the particular empirical normalization used in that computation.

\section{Constructibility Status}

Let
\[
  c_N=2\sin\!\left(\frac{\pi}{N}\right),
  \qquad P_N=\frac{N}{2}c_N.
\]
For the dyadic branch, starting from $c_6=1$, the stable bisection recurrence is
\[
  c_{2N}
  =\frac{c_N}{\sqrt{2+\sqrt{4-c_N^2}}}.
\]
It uses only rational operations and square roots.  Therefore all dyadic
polygon values and all rationally weighted dyadic accelerants are
compass-and-straightedge constructible.

For exact tripling, the triple-angle identity gives
\[
  c_N=3c_{3N}-c_{3N}^3.
\]
Hence $u=c_{3N}$ is the small positive root of
\[
  u^3-3u+c_N=0.
\]
This cubic is not generally solvable by compass and straightedge.  In
particular, the regular $18$-gon is not constructible, so
\[
  6\to18\to54\to\cdots
\]
cannot be presented as a ladder of exact constructible regular polygons.

\subsection{A Constructible Linear Seed}

\begin{proposition}[Linear tripling seed]
Define
\[
  \widetilde c^{(0)}_{3N}=\frac{c_N}{3}.
\]
This segment is constructible from $c_N$, and
\[
  \widetilde c^{(0)}_{3N}-c_{3N}
  =-\frac{8\pi^3}{81}N^{-3}+O(N^{-5}).
\]
In particular, the chord seed error is $O(N^{-3})$.
\end{proposition}

\begin{proof}
Set $x=\pi/N$.  Then
\[
  \frac{c_N}{3}=\frac{2}{3}\sin x
  =\frac{2x}{3}-\frac{x^3}{9}+O(x^5),
\]
whereas
\[
  c_{3N}=2\sin(x/3)
  =\frac{2x}{3}-\frac{x^3}{81}+O(x^5).
\]
Subtract and replace $x$ by $\pi/N$.
\end{proof}

\subsection{Perturbation Requirement}

Suppose the chord used at scale $M=b^mN$ is perturbed by $\delta_M$.  Since
$P_M=(M/2)c_M$, the accelerated perturbation is bounded by
\[
  \frac{N}{2}\sum_{m=0}^{r-1}
  |w_m^{(r,b)}|b^m|\delta_{b^mN}|.
\]

\begin{theorem}[Sufficient perturbation exponent]
For fixed $b$ and $r$, suppose
\[
  |\delta_M|\le AM^{-\alpha}
\]
for all scales used by the operator.  If
\[
  \alpha\ge2r+1,
\]
then the perturbation contribution is $O(N^{-2r})$ and therefore does not
change the formal order of the accelerated hierarchy.
\end{theorem}

\begin{proof}
Substitution gives
\[
\frac{N}{2}\sum_m|w_m|b^m|\delta_{b^mN}|
\le
\frac{A}{2}N^{1-\alpha}
\sum_m|w_m|b^{m(1-\alpha)}.
\]
For fixed $r$, the finite weighted sum is independent of $N$.  If
$\alpha\ge2r+1$, then $N^{1-\alpha}=O(N^{-2r})$.
\end{proof}

\begin{remark}[Conditional cubic refinement]
If an independently documented constructible refinement map has a local error
law
\[
  e_+=Ce^3+O(e^5)
\]
and if the $O(N^{-3})$ seed lies in its uniform local basin, then after $s$
refinements its error exponent becomes $3^{s+1}$.  The exponent condition
\[
  3^{s+1}\ge2r+1
\]
is sufficient asymptotically.  This statement is conditional on the actual
geometric realization and its domain; the present paper does not treat odd
symmetry alone as proof that the linear term vanishes, nor does it claim a
standalone exact trisection operator.
\end{remark}

\section{Limitations}

The fixed-$r$ acceleration law is polynomial in the resolution variable and
is not a competitor to AGM or modular-form algorithms as an arithmetic method
for computing $\pi$.  The numerical branch comparison depends on the selected
cost normalization.  Exact tripling data generated by cubic root-solving are
valid for testing the scale architecture but are not straightedge-and-compass
constructions.  A fully geometric tripling implementation requires a separate,
explicit construction and an error certificate satisfying the perturbation
condition above.

The paper also does not claim that the cancellation algebra is new.  It is a
Richardson-type extrapolation.  The contribution is the geometric-node
stability analysis, the explicit remainder control, the scale/work comparison,
and the precise separation of exact dyadic constructibility from approximate
tripling.

\section{Conclusion}

The polygonal sequence is a scale-indexed carrier of removable error
structure.  The scale-$b$ operator cancels the first $r-1$ even-power modes,
its geometric-node weights are uniformly stable, and the remaining sine tail
has an explicit factorial bound.  Larger scales improve weight stability but
also increase work, so optimization requires a cost model rather than a
largest-scale rule.

The constructibility conclusion is equally precise.  Dyadic refinement is
exactly constructible.  Exact tripling is generally not.  Approximate tripling
can nevertheless be studied quantitatively: the linear chord seed is
constructible with $O(N^{-3})$ error, and any refinement mechanism meeting the
required perturbation exponent preserves the accelerated order.

\[
  \boxed{\text{The circle is approached not only by refinement, but by
  cancellation of refinement error.}}
\]

\begin{thebibliography}{9}

\bibitem{archimedes}
Archimedes,
\emph{Measurement of a Circle}, in T. L. Heath, ed.,
\emph{The Works of Archimedes}, Cambridge University Press, 1897.

\bibitem{richardson}
L. F. Richardson,
``The approximate arithmetical solution by finite differences of physical
problems involving differential equations, with an application to the
stresses in a masonry dam,''
\emph{Philosophical Transactions of the Royal Society A}, 210 (1911),
307--357.

\bibitem{trefethen}
L. N. Trefethen and J. A. C. Weideman,
``The exponentially convergent trapezoidal rule,''
\emph{SIAM Review}, 56(3) (2014), 385--458.

\bibitem{borwein}
J. M. Borwein and P. B. Borwein,
\emph{Pi and the AGM: A Study in Analytic Number Theory and Computational
Complexity}, Wiley, 1987.

\end{thebibliography}

\end{document}
